% Ngân hàng câu hỏi TOÁN THỰC TẾ T10 -- Chương 6, Bài 16, Dạng D72
%% Dạng: Hàm số bậc hai trong thực tiễn: ném vật, lợi nhuận
%% 100 câu: NB 30 -- TH 40 -- VD 20 -- VDC 10 ; mỗi mức: trắc nghiệm : đúng/sai : trả lời ngắn = 3 : 1 : 1
%% Phương án ĐẦU TIÊN trong \cauTN luôn là đáp án đúng (gói trondethi tự đảo).
%% Phần II: \y{nội dung}{d|s}. Phần III: đáp án đặt ở đối số cuối.

\cauTN[id=T10-C6-B16-D72-NB-001,muc=NB,loigiai={Đồ thị là parabol có $a=-4{,}9<0$ nên $h$ đạt giá trị lớn nhất tại $t=-\dfrac{b}{2a}=\dfrac{14{,}7}{9{,}8}=1{,}5$ giây. Khi đó $h=12{,}02$ mét.}]{Một quả bóng được đá lên từ độ cao $1$ m với vận tốc ban đầu theo phương thẳng đứng là $14{,}7$ m/s. Độ cao của bóng sau $t$ giây là $h(t)=-4{,}9t^2+14{,}7t+1$ (mét). Độ cao lớn nhất của quả bóng là bao nhiêu mét (làm tròn đến hàng phần trăm)?}
  {$12{,}02$.}
  {$13{,}22$.}{$11{,}22$.}{$14{,}7$.}

\cauTN[id=T10-C6-B16-D72-NB-002,muc=NB,loigiai={Đồ thị là parabol có $a=-4{,}9<0$ nên $h$ đạt giá trị lớn nhất tại $t=-\dfrac{b}{2a}=\dfrac{19{,}6}{9{,}8}=2$ giây. Khi đó $h=19{,}6$ mét.}]{Một quả bóng được đá lên từ độ cao $0$ m với vận tốc ban đầu theo phương thẳng đứng là $19{,}6$ m/s. Độ cao của bóng sau $t$ giây là $h(t)=-4{,}9t^2+19{,}6t+0$ (mét). Độ cao lớn nhất của quả bóng là bao nhiêu mét (làm tròn đến hàng phần trăm)?}
  {$19{,}6$.}
  {$20{,}8$.}{$18{,}8$.}{$22{,}2$.}

\cauTN[id=T10-C6-B16-D72-NB-003,muc=NB,loigiai={Đồ thị là parabol có $a=-4{,}9<0$ nên $h$ đạt giá trị lớn nhất tại $t=-\dfrac{b}{2a}=\dfrac{14{,}7}{9{,}8}=1{,}5$ giây. Khi đó $h=12{,}02$ mét.}]{Một quả bóng được đá lên từ độ cao $1$ m với vận tốc ban đầu theo phương thẳng đứng là $14{,}7$ m/s. Độ cao của bóng sau $t$ giây là $h(t)=-4{,}9t^2+14{,}7t+1$ (mét). Độ cao lớn nhất của quả bóng là bao nhiêu mét (làm tròn đến hàng phần trăm)?}
  {$12{,}02$.}
  {$13{,}22$.}{$11{,}22$.}{$14{,}7$.}

\cauTN[id=T10-C6-B16-D72-NB-004,muc=NB,loigiai={Đồ thị là parabol có $a=-4{,}9<0$ nên $h$ đạt giá trị lớn nhất tại $t=-\dfrac{b}{2a}=\dfrac{19{,}6}{9{,}8}=2$ giây. Khi đó $h=21{,}1$ mét.}]{Một quả bóng được đá lên từ độ cao $1{,}5$ m với vận tốc ban đầu theo phương thẳng đứng là $19{,}6$ m/s. Độ cao của bóng sau $t$ giây là $h(t)=-4{,}9t^2+19{,}6t+1{,}5$ (mét). Độ cao lớn nhất của quả bóng là bao nhiêu mét (làm tròn đến hàng phần trăm)?}
  {$21{,}1$.}
  {$22{,}3$.}{$20{,}3$.}{$19{,}6$.}

\cauTN[id=T10-C6-B16-D72-NB-005,muc=NB,loigiai={Đồ thị là parabol có $a=-4{,}9<0$ nên $h$ đạt giá trị lớn nhất tại $t=-\dfrac{b}{2a}=\dfrac{14{,}7}{9{,}8}=1{,}5$ giây. Khi đó $h=13{,}02$ mét.}]{Một quả bóng được đá lên từ độ cao $2$ m với vận tốc ban đầu theo phương thẳng đứng là $14{,}7$ m/s. Độ cao của bóng sau $t$ giây là $h(t)=-4{,}9t^2+14{,}7t+2$ (mét). Độ cao lớn nhất của quả bóng là bao nhiêu mét (làm tròn đến hàng phần trăm)?}
  {$13{,}02$.}
  {$14{,}22$.}{$12{,}22$.}{$14{,}7$.}

\cauTN[id=T10-C6-B16-D72-NB-006,muc=NB,loigiai={Đồ thị là parabol có $a=-4{,}9<0$ nên $h$ đạt giá trị lớn nhất tại $t=-\dfrac{b}{2a}=\dfrac{14{,}7}{9{,}8}=1{,}5$ giây. Khi đó $h=13{,}02$ mét.}]{Một quả bóng được đá lên từ độ cao $2$ m với vận tốc ban đầu theo phương thẳng đứng là $14{,}7$ m/s. Độ cao của bóng sau $t$ giây là $h(t)=-4{,}9t^2+14{,}7t+2$ (mét). Độ cao lớn nhất của quả bóng là bao nhiêu mét (làm tròn đến hàng phần trăm)?}
  {$13{,}02$.}
  {$14{,}22$.}{$12{,}22$.}{$14{,}7$.}

\cauTN[id=T10-C6-B16-D72-NB-007,muc=NB,loigiai={Đồ thị là parabol có $a=-4{,}9<0$ nên $h$ đạt giá trị lớn nhất tại $t=-\dfrac{b}{2a}=\dfrac{29{,}4}{9{,}8}=3$ giây. Khi đó $h=46{,}1$ mét.}]{Một quả bóng được đá lên từ độ cao $2$ m với vận tốc ban đầu theo phương thẳng đứng là $29{,}4$ m/s. Độ cao của bóng sau $t$ giây là $h(t)=-4{,}9t^2+29{,}4t+2$ (mét). Độ cao lớn nhất của quả bóng là bao nhiêu mét (làm tròn đến hàng phần trăm)?}
  {$46{,}1$.}
  {$47{,}3$.}{$45{,}3$.}{$29{,}4$.}

\cauTN[id=T10-C6-B16-D72-NB-008,muc=NB,loigiai={Đồ thị là parabol có $a=-4{,}9<0$ nên $h$ đạt giá trị lớn nhất tại $t=-\dfrac{b}{2a}=\dfrac{14{,}7}{9{,}8}=1{,}5$ giây. Khi đó $h=13{,}02$ mét.}]{Một quả bóng được đá lên từ độ cao $2$ m với vận tốc ban đầu theo phương thẳng đứng là $14{,}7$ m/s. Độ cao của bóng sau $t$ giây là $h(t)=-4{,}9t^2+14{,}7t+2$ (mét). Độ cao lớn nhất của quả bóng là bao nhiêu mét (làm tròn đến hàng phần trăm)?}
  {$13{,}02$.}
  {$14{,}22$.}{$12{,}22$.}{$14{,}7$.}

\cauTN[id=T10-C6-B16-D72-NB-009,muc=NB,loigiai={Đồ thị là parabol có $a=-4{,}9<0$ nên $h$ đạt giá trị lớn nhất tại $t=-\dfrac{b}{2a}=\dfrac{19{,}6}{9{,}8}=2$ giây. Khi đó $h=19{,}6$ mét.}]{Một quả bóng được đá lên từ độ cao $0$ m với vận tốc ban đầu theo phương thẳng đứng là $19{,}6$ m/s. Độ cao của bóng sau $t$ giây là $h(t)=-4{,}9t^2+19{,}6t+0$ (mét). Độ cao lớn nhất của quả bóng là bao nhiêu mét (làm tròn đến hàng phần trăm)?}
  {$19{,}6$.}
  {$20{,}8$.}{$18{,}8$.}{$22{,}2$.}

\cauTN[id=T10-C6-B16-D72-NB-010,muc=NB,loigiai={Đồ thị là parabol có $a=-4{,}9<0$ nên $h$ đạt giá trị lớn nhất tại $t=-\dfrac{b}{2a}=\dfrac{19{,}6}{9{,}8}=2$ giây. Khi đó $h=20{,}6$ mét.}]{Một quả bóng được đá lên từ độ cao $1$ m với vận tốc ban đầu theo phương thẳng đứng là $19{,}6$ m/s. Độ cao của bóng sau $t$ giây là $h(t)=-4{,}9t^2+19{,}6t+1$ (mét). Độ cao lớn nhất của quả bóng là bao nhiêu mét (làm tròn đến hàng phần trăm)?}
  {$20{,}6$.}
  {$21{,}8$.}{$19{,}8$.}{$19{,}6$.}

\cauTN[id=T10-C6-B16-D72-NB-011,muc=NB,loigiai={Đồ thị là parabol có $a=-4{,}9<0$ nên $h$ đạt giá trị lớn nhất tại $t=-\dfrac{b}{2a}=\dfrac{24{,}5}{9{,}8}=2{,}5$ giây. Khi đó $h=31{,}62$ mét.}]{Một quả bóng được đá lên từ độ cao $1$ m với vận tốc ban đầu theo phương thẳng đứng là $24{,}5$ m/s. Độ cao của bóng sau $t$ giây là $h(t)=-4{,}9t^2+24{,}5t+1$ (mét). Độ cao lớn nhất của quả bóng là bao nhiêu mét (làm tròn đến hàng phần trăm)?}
  {$31{,}62$.}
  {$32{,}83$.}{$30{,}82$.}{$24{,}5$.}

\cauTN[id=T10-C6-B16-D72-NB-012,muc=NB,loigiai={Đồ thị là parabol có $a=-4{,}9<0$ nên $h$ đạt giá trị lớn nhất tại $t=-\dfrac{b}{2a}=\dfrac{14{,}7}{9{,}8}=1{,}5$ giây. Khi đó $h=12{,}02$ mét.}]{Một quả bóng được đá lên từ độ cao $1$ m với vận tốc ban đầu theo phương thẳng đứng là $14{,}7$ m/s. Độ cao của bóng sau $t$ giây là $h(t)=-4{,}9t^2+14{,}7t+1$ (mét). Độ cao lớn nhất của quả bóng là bao nhiêu mét (làm tròn đến hàng phần trăm)?}
  {$12{,}02$.}
  {$13{,}22$.}{$11{,}22$.}{$14{,}7$.}

\cauTN[id=T10-C6-B16-D72-NB-013,muc=NB,loigiai={Đồ thị là parabol có $a=-4{,}9<0$ nên $h$ đạt giá trị lớn nhất tại $t=-\dfrac{b}{2a}=\dfrac{29{,}4}{9{,}8}=3$ giây. Khi đó $h=46{,}1$ mét.}]{Một quả bóng được đá lên từ độ cao $2$ m với vận tốc ban đầu theo phương thẳng đứng là $29{,}4$ m/s. Độ cao của bóng sau $t$ giây là $h(t)=-4{,}9t^2+29{,}4t+2$ (mét). Độ cao lớn nhất của quả bóng là bao nhiêu mét (làm tròn đến hàng phần trăm)?}
  {$46{,}1$.}
  {$47{,}3$.}{$45{,}3$.}{$29{,}4$.}

\cauTN[id=T10-C6-B16-D72-NB-014,muc=NB,loigiai={Đồ thị là parabol có $a=-4{,}9<0$ nên $h$ đạt giá trị lớn nhất tại $t=-\dfrac{b}{2a}=\dfrac{29{,}4}{9{,}8}=3$ giây. Khi đó $h=45{,}6$ mét.}]{Một quả bóng được đá lên từ độ cao $1{,}5$ m với vận tốc ban đầu theo phương thẳng đứng là $29{,}4$ m/s. Độ cao của bóng sau $t$ giây là $h(t)=-4{,}9t^2+29{,}4t+1{,}5$ (mét). Độ cao lớn nhất của quả bóng là bao nhiêu mét (làm tròn đến hàng phần trăm)?}
  {$45{,}6$.}
  {$46{,}8$.}{$44{,}8$.}{$29{,}4$.}

\cauTN[id=T10-C6-B16-D72-NB-015,muc=NB,loigiai={Đồ thị là parabol có $a=-4{,}9<0$ nên $h$ đạt giá trị lớn nhất tại $t=-\dfrac{b}{2a}=\dfrac{19{,}6}{9{,}8}=2$ giây. Khi đó $h=21{,}6$ mét.}]{Một quả bóng được đá lên từ độ cao $2$ m với vận tốc ban đầu theo phương thẳng đứng là $19{,}6$ m/s. Độ cao của bóng sau $t$ giây là $h(t)=-4{,}9t^2+19{,}6t+2$ (mét). Độ cao lớn nhất của quả bóng là bao nhiêu mét (làm tròn đến hàng phần trăm)?}
  {$21{,}6$.}
  {$22{,}8$.}{$20{,}8$.}{$19{,}6$.}

\cauTN[id=T10-C6-B16-D72-NB-016,muc=NB,loigiai={Đồ thị là parabol có $a=-4{,}9<0$ nên $h$ đạt giá trị lớn nhất tại $t=-\dfrac{b}{2a}=\dfrac{29{,}4}{9{,}8}=3$ giây. Khi đó $h=45{,}1$ mét.}]{Một quả bóng được đá lên từ độ cao $1$ m với vận tốc ban đầu theo phương thẳng đứng là $29{,}4$ m/s. Độ cao của bóng sau $t$ giây là $h(t)=-4{,}9t^2+29{,}4t+1$ (mét). Độ cao lớn nhất của quả bóng là bao nhiêu mét (làm tròn đến hàng phần trăm)?}
  {$45{,}1$.}
  {$46{,}3$.}{$44{,}3$.}{$29{,}4$.}

\cauTN[id=T10-C6-B16-D72-NB-017,muc=NB,loigiai={Đồ thị là parabol có $a=-4{,}9<0$ nên $h$ đạt giá trị lớn nhất tại $t=-\dfrac{b}{2a}=\dfrac{29{,}4}{9{,}8}=3$ giây. Khi đó $h=45{,}1$ mét.}]{Một quả bóng được đá lên từ độ cao $1$ m với vận tốc ban đầu theo phương thẳng đứng là $29{,}4$ m/s. Độ cao của bóng sau $t$ giây là $h(t)=-4{,}9t^2+29{,}4t+1$ (mét). Độ cao lớn nhất của quả bóng là bao nhiêu mét (làm tròn đến hàng phần trăm)?}
  {$45{,}1$.}
  {$46{,}3$.}{$44{,}3$.}{$29{,}4$.}

\cauTN[id=T10-C6-B16-D72-NB-018,muc=NB,loigiai={Đồ thị là parabol có $a=-4{,}9<0$ nên $h$ đạt giá trị lớn nhất tại $t=-\dfrac{b}{2a}=\dfrac{14{,}7}{9{,}8}=1{,}5$ giây. Khi đó $h=12{,}02$ mét.}]{Một quả bóng được đá lên từ độ cao $1$ m với vận tốc ban đầu theo phương thẳng đứng là $14{,}7$ m/s. Độ cao của bóng sau $t$ giây là $h(t)=-4{,}9t^2+14{,}7t+1$ (mét). Độ cao lớn nhất của quả bóng là bao nhiêu mét (làm tròn đến hàng phần trăm)?}
  {$12{,}02$.}
  {$13{,}22$.}{$11{,}22$.}{$14{,}7$.}

\cauDS[id=T10-C6-B16-D72-NB-019,muc=NB,loigiai={Đồ thị là parabol có $a=-4{,}9<0$ nên $h$ đạt giá trị lớn nhất tại $t=-\dfrac{b}{2a}=\dfrac{19{,}6}{9{,}8}=2$ giây. Khi đó $h=20{,}6$ mét.}]{Một quả bóng được đá lên từ độ cao $1$ m với vận tốc ban đầu theo phương thẳng đứng là $19{,}6$ m/s. Độ cao của bóng sau $t$ giây là $h(t)=-4{,}9t^2+19{,}6t+1$ (mét).}{%
  \y{Hệ số $a$ của hàm số là $-4{,}9$.}{d}
  \y{Bóng đạt độ cao lớn nhất tại thời điểm $t=2$ giây.}{d}
  \y{Độ cao ban đầu của bóng là $1$ mét.}{d}
  \y{Độ cao lớn nhất của bóng là $22{,}6$ mét.}{s}}

\cauDS[id=T10-C6-B16-D72-NB-020,muc=NB,loigiai={Đồ thị là parabol có $a=-4{,}9<0$ nên $h$ đạt giá trị lớn nhất tại $t=-\dfrac{b}{2a}=\dfrac{24{,}5}{9{,}8}=2{,}5$ giây. Khi đó $h=31{,}62$ mét.}]{Một quả bóng được đá lên từ độ cao $1$ m với vận tốc ban đầu theo phương thẳng đứng là $24{,}5$ m/s. Độ cao của bóng sau $t$ giây là $h(t)=-4{,}9t^2+24{,}5t+1$ (mét).}{%
  \y{Hệ số $a$ của hàm số là $-4{,}9$.}{d}
  \y{Bóng đạt độ cao lớn nhất tại thời điểm $t=2{,}5$ giây.}{d}
  \y{Độ cao ban đầu của bóng là $1$ mét.}{d}
  \y{Độ cao lớn nhất của bóng là $33{,}62$ mét.}{s}}

\cauDS[id=T10-C6-B16-D72-NB-021,muc=NB,loigiai={Đồ thị là parabol có $a=-4{,}9<0$ nên $h$ đạt giá trị lớn nhất tại $t=-\dfrac{b}{2a}=\dfrac{29{,}4}{9{,}8}=3$ giây. Khi đó $h=45{,}1$ mét.}]{Một quả bóng được đá lên từ độ cao $1$ m với vận tốc ban đầu theo phương thẳng đứng là $29{,}4$ m/s. Độ cao của bóng sau $t$ giây là $h(t)=-4{,}9t^2+29{,}4t+1$ (mét).}{%
  \y{Hệ số $a$ của hàm số là $-4{,}9$.}{d}
  \y{Bóng đạt độ cao lớn nhất tại thời điểm $t=3$ giây.}{d}
  \y{Độ cao ban đầu của bóng là $1$ mét.}{d}
  \y{Độ cao lớn nhất của bóng là $47{,}1$ mét.}{s}}

\cauDS[id=T10-C6-B16-D72-NB-022,muc=NB,loigiai={Đồ thị là parabol có $a=-4{,}9<0$ nên $h$ đạt giá trị lớn nhất tại $t=-\dfrac{b}{2a}=\dfrac{19{,}6}{9{,}8}=2$ giây. Khi đó $h=20{,}6$ mét.}]{Một quả bóng được đá lên từ độ cao $1$ m với vận tốc ban đầu theo phương thẳng đứng là $19{,}6$ m/s. Độ cao của bóng sau $t$ giây là $h(t)=-4{,}9t^2+19{,}6t+1$ (mét).}{%
  \y{Hệ số $a$ của hàm số là $-4{,}9$.}{d}
  \y{Bóng đạt độ cao lớn nhất tại thời điểm $t=2$ giây.}{d}
  \y{Độ cao ban đầu của bóng là $1$ mét.}{d}
  \y{Độ cao lớn nhất của bóng là $22{,}6$ mét.}{s}}

\cauDS[id=T10-C6-B16-D72-NB-023,muc=NB,loigiai={Đồ thị là parabol có $a=-4{,}9<0$ nên $h$ đạt giá trị lớn nhất tại $t=-\dfrac{b}{2a}=\dfrac{29{,}4}{9{,}8}=3$ giây. Khi đó $h=45{,}6$ mét.}]{Một quả bóng được đá lên từ độ cao $1{,}5$ m với vận tốc ban đầu theo phương thẳng đứng là $29{,}4$ m/s. Độ cao của bóng sau $t$ giây là $h(t)=-4{,}9t^2+29{,}4t+1{,}5$ (mét).}{%
  \y{Hệ số $a$ của hàm số là $-4{,}9$.}{d}
  \y{Bóng đạt độ cao lớn nhất tại thời điểm $t=3$ giây.}{d}
  \y{Độ cao ban đầu của bóng là $1{,}5$ mét.}{d}
  \y{Độ cao lớn nhất của bóng là $47{,}6$ mét.}{s}}

\cauDS[id=T10-C6-B16-D72-NB-024,muc=NB,loigiai={Đồ thị là parabol có $a=-4{,}9<0$ nên $h$ đạt giá trị lớn nhất tại $t=-\dfrac{b}{2a}=\dfrac{19{,}6}{9{,}8}=2$ giây. Khi đó $h=21{,}1$ mét.}]{Một quả bóng được đá lên từ độ cao $1{,}5$ m với vận tốc ban đầu theo phương thẳng đứng là $19{,}6$ m/s. Độ cao của bóng sau $t$ giây là $h(t)=-4{,}9t^2+19{,}6t+1{,}5$ (mét).}{%
  \y{Hệ số $a$ của hàm số là $-4{,}9$.}{d}
  \y{Bóng đạt độ cao lớn nhất tại thời điểm $t=2$ giây.}{d}
  \y{Độ cao ban đầu của bóng là $1{,}5$ mét.}{d}
  \y{Độ cao lớn nhất của bóng là $23{,}1$ mét.}{s}}

\cauTLN[id=T10-C6-B16-D72-NB-025,muc=NB,loigiai={Đồ thị là parabol có $a=-4{,}9<0$ nên $h$ đạt giá trị lớn nhất tại $t=-\dfrac{b}{2a}=\dfrac{19{,}6}{9{,}8}=2$ giây. Khi đó $h=19{,}6$ mét.}]{Một quả bóng được đá lên từ độ cao $0$ m với vận tốc ban đầu theo phương thẳng đứng là $19{,}6$ m/s. Độ cao của bóng sau $t$ giây là $h(t)=-4{,}9t^2+19{,}6t+0$ (mét). Độ cao lớn nhất của quả bóng là bao nhiêu mét (làm tròn đến hàng phần trăm)?}{19,6}

\cauTLN[id=T10-C6-B16-D72-NB-026,muc=NB,loigiai={Đồ thị là parabol có $a=-4{,}9<0$ nên $h$ đạt giá trị lớn nhất tại $t=-\dfrac{b}{2a}=\dfrac{29{,}4}{9{,}8}=3$ giây. Khi đó $h=45{,}1$ mét.}]{Một quả bóng được đá lên từ độ cao $1$ m với vận tốc ban đầu theo phương thẳng đứng là $29{,}4$ m/s. Độ cao của bóng sau $t$ giây là $h(t)=-4{,}9t^2+29{,}4t+1$ (mét). Độ cao lớn nhất của quả bóng là bao nhiêu mét (làm tròn đến hàng phần trăm)?}{45,1}

\cauTLN[id=T10-C6-B16-D72-NB-027,muc=NB,loigiai={Đồ thị là parabol có $a=-4{,}9<0$ nên $h$ đạt giá trị lớn nhất tại $t=-\dfrac{b}{2a}=\dfrac{29{,}4}{9{,}8}=3$ giây. Khi đó $h=45{,}6$ mét.}]{Một quả bóng được đá lên từ độ cao $1{,}5$ m với vận tốc ban đầu theo phương thẳng đứng là $29{,}4$ m/s. Độ cao của bóng sau $t$ giây là $h(t)=-4{,}9t^2+29{,}4t+1{,}5$ (mét). Độ cao lớn nhất của quả bóng là bao nhiêu mét (làm tròn đến hàng phần trăm)?}{45,6}

\cauTLN[id=T10-C6-B16-D72-NB-028,muc=NB,loigiai={Đồ thị là parabol có $a=-4{,}9<0$ nên $h$ đạt giá trị lớn nhất tại $t=-\dfrac{b}{2a}=\dfrac{29{,}4}{9{,}8}=3$ giây. Khi đó $h=45{,}1$ mét.}]{Một quả bóng được đá lên từ độ cao $1$ m với vận tốc ban đầu theo phương thẳng đứng là $29{,}4$ m/s. Độ cao của bóng sau $t$ giây là $h(t)=-4{,}9t^2+29{,}4t+1$ (mét). Độ cao lớn nhất của quả bóng là bao nhiêu mét (làm tròn đến hàng phần trăm)?}{45,1}

\cauTLN[id=T10-C6-B16-D72-NB-029,muc=NB,loigiai={Đồ thị là parabol có $a=-4{,}9<0$ nên $h$ đạt giá trị lớn nhất tại $t=-\dfrac{b}{2a}=\dfrac{24{,}5}{9{,}8}=2{,}5$ giây. Khi đó $h=30{,}62$ mét.}]{Một quả bóng được đá lên từ độ cao $0$ m với vận tốc ban đầu theo phương thẳng đứng là $24{,}5$ m/s. Độ cao của bóng sau $t$ giây là $h(t)=-4{,}9t^2+24{,}5t+0$ (mét). Độ cao lớn nhất của quả bóng là bao nhiêu mét (làm tròn đến hàng phần trăm)?}{30,62}

\cauTLN[id=T10-C6-B16-D72-NB-030,muc=NB,loigiai={Đồ thị là parabol có $a=-4{,}9<0$ nên $h$ đạt giá trị lớn nhất tại $t=-\dfrac{b}{2a}=\dfrac{29{,}4}{9{,}8}=3$ giây. Khi đó $h=46{,}1$ mét.}]{Một quả bóng được đá lên từ độ cao $2$ m với vận tốc ban đầu theo phương thẳng đứng là $29{,}4$ m/s. Độ cao của bóng sau $t$ giây là $h(t)=-4{,}9t^2+29{,}4t+2$ (mét). Độ cao lớn nhất của quả bóng là bao nhiêu mét (làm tròn đến hàng phần trăm)?}{46,1}

\cauTN[id=T10-C6-B16-D72-TH-001,muc=TH,loigiai={Đồ thị là parabol có $a=-4{,}9<0$ nên $h$ đạt giá trị lớn nhất tại $t=-\dfrac{b}{2a}=\dfrac{24{,}5}{9{,}8}=2{,}5$ giây. Khi đó $h=30{,}62$ mét.}]{Một quả bóng được đá lên từ độ cao $0$ m với vận tốc ban đầu theo phương thẳng đứng là $24{,}5$ m/s. Độ cao của bóng sau $t$ giây là $h(t)=-4{,}9t^2+24{,}5t+0$ (mét). Độ cao lớn nhất của quả bóng là bao nhiêu mét (làm tròn đến hàng phần trăm)?}
  {$30{,}62$.}
  {$31{,}82$.}{$29{,}82$.}{$24{,}5$.}

\cauTN[id=T10-C6-B16-D72-TH-002,muc=TH,loigiai={$a=-1<0$ nên $L$ lớn nhất tại $x=\dfrac{86}{2\cdot1}=43$, khi đó $L=1\,955$ nghìn đồng.}]{Một cửa hàng ước tính lợi nhuận trong ngày thứ $x$ của tháng cho bởi $L(x)=-x^2+86x+106$ (nghìn đồng), với $1\le x\le 30$. Lợi nhuận lớn nhất trong tháng là bao nhiêu nghìn đồng?}
  {$1\,955$.}
  {$2\,005$.}{$106$.}{$1\,925$.}

\cauTN[id=T10-C6-B16-D72-TH-003,muc=TH,loigiai={Đồ thị là parabol có $a=-4{,}9<0$ nên $h$ đạt giá trị lớn nhất tại $t=-\dfrac{b}{2a}=\dfrac{29{,}4}{9{,}8}=3$ giây. Khi đó $h=44{,}1$ mét.}]{Một quả bóng được đá lên từ độ cao $0$ m với vận tốc ban đầu theo phương thẳng đứng là $29{,}4$ m/s. Độ cao của bóng sau $t$ giây là $h(t)=-4{,}9t^2+29{,}4t+0$ (mét). Độ cao lớn nhất của quả bóng là bao nhiêu mét (làm tròn đến hàng phần trăm)?}
  {$44{,}1$.}
  {$45{,}3$.}{$43{,}3$.}{$29{,}4$.}

\cauTN[id=T10-C6-B16-D72-TH-004,muc=TH,loigiai={$a=-3<0$ nên $L$ lớn nhất tại $x=\dfrac{81}{2\cdot3}=\dfrac{27}{2}$, khi đó $L=1347{,}75$ nghìn đồng.}]{Một cửa hàng ước tính lợi nhuận trong ngày thứ $x$ của tháng cho bởi $L(x)=-3x^2+81x+801$ (nghìn đồng), với $1\le x\le 30$. Lợi nhuận lớn nhất trong tháng là bao nhiêu nghìn đồng?}
  {$1347{,}75$.}
  {$1397{,}75$.}{$801$.}{$1317{,}75$.}

\cauTN[id=T10-C6-B16-D72-TH-005,muc=TH,loigiai={Đồ thị là parabol có $a=-4{,}9<0$ nên $h$ đạt giá trị lớn nhất tại $t=-\dfrac{b}{2a}=\dfrac{19{,}6}{9{,}8}=2$ giây. Khi đó $h=19{,}6$ mét.}]{Một quả bóng được đá lên từ độ cao $0$ m với vận tốc ban đầu theo phương thẳng đứng là $19{,}6$ m/s. Độ cao của bóng sau $t$ giây là $h(t)=-4{,}9t^2+19{,}6t+0$ (mét). Độ cao lớn nhất của quả bóng là bao nhiêu mét (làm tròn đến hàng phần trăm)?}
  {$19{,}6$.}
  {$20{,}8$.}{$18{,}8$.}{$22{,}2$.}

\cauTN[id=T10-C6-B16-D72-TH-006,muc=TH,loigiai={$a=-3<0$ nên $L$ lớn nhất tại $x=\dfrac{80}{2\cdot3}=\dfrac{40}{3}$, khi đó $L=764{,}33$ nghìn đồng.}]{Một cửa hàng ước tính lợi nhuận trong ngày thứ $x$ của tháng cho bởi $L(x)=-3x^2+80x+231$ (nghìn đồng), với $1\le x\le 30$. Lợi nhuận lớn nhất trong tháng là bao nhiêu nghìn đồng?}
  {$764{,}33$.}
  {$814{,}33$.}{$231$.}{$734{,}33$.}

\cauTN[id=T10-C6-B16-D72-TH-007,muc=TH,loigiai={Đồ thị là parabol có $a=-4{,}9<0$ nên $h$ đạt giá trị lớn nhất tại $t=-\dfrac{b}{2a}=\dfrac{19{,}6}{9{,}8}=2$ giây. Khi đó $h=21{,}1$ mét.}]{Một quả bóng được đá lên từ độ cao $1{,}5$ m với vận tốc ban đầu theo phương thẳng đứng là $19{,}6$ m/s. Độ cao của bóng sau $t$ giây là $h(t)=-4{,}9t^2+19{,}6t+1{,}5$ (mét). Độ cao lớn nhất của quả bóng là bao nhiêu mét (làm tròn đến hàng phần trăm)?}
  {$21{,}1$.}
  {$22{,}3$.}{$20{,}3$.}{$19{,}6$.}

\cauTN[id=T10-C6-B16-D72-TH-008,muc=TH,loigiai={$a=-1<0$ nên $L$ lớn nhất tại $x=\dfrac{38}{2\cdot1}=19$, khi đó $L=705$ nghìn đồng.}]{Một cửa hàng ước tính lợi nhuận trong ngày thứ $x$ của tháng cho bởi $L(x)=-x^2+38x+344$ (nghìn đồng), với $1\le x\le 30$. Lợi nhuận lớn nhất trong tháng là bao nhiêu nghìn đồng?}
  {$705$.}
  {$755$.}{$344$.}{$675$.}

\cauTN[id=T10-C6-B16-D72-TH-009,muc=TH,loigiai={Đồ thị là parabol có $a=-4{,}9<0$ nên $h$ đạt giá trị lớn nhất tại $t=-\dfrac{b}{2a}=\dfrac{24{,}5}{9{,}8}=2{,}5$ giây. Khi đó $h=32{,}12$ mét.}]{Một quả bóng được đá lên từ độ cao $1{,}5$ m với vận tốc ban đầu theo phương thẳng đứng là $24{,}5$ m/s. Độ cao của bóng sau $t$ giây là $h(t)=-4{,}9t^2+24{,}5t+1{,}5$ (mét). Độ cao lớn nhất của quả bóng là bao nhiêu mét (làm tròn đến hàng phần trăm)?}
  {$32{,}12$.}
  {$33{,}33$.}{$31{,}32$.}{$24{,}5$.}

\cauTN[id=T10-C6-B16-D72-TH-010,muc=TH,loigiai={$a=-4<0$ nên $L$ lớn nhất tại $x=\dfrac{69}{2\cdot4}=\dfrac{69}{8}$, khi đó $L=586{,}56$ nghìn đồng.}]{Một cửa hàng ước tính lợi nhuận trong ngày thứ $x$ của tháng cho bởi $L(x)=-4x^2+69x+289$ (nghìn đồng), với $1\le x\le 30$. Lợi nhuận lớn nhất trong tháng là bao nhiêu nghìn đồng?}
  {$586{,}56$.}
  {$636{,}56$.}{$289$.}{$556{,}56$.}

\cauTN[id=T10-C6-B16-D72-TH-011,muc=TH,loigiai={Đồ thị là parabol có $a=-4{,}9<0$ nên $h$ đạt giá trị lớn nhất tại $t=-\dfrac{b}{2a}=\dfrac{29{,}4}{9{,}8}=3$ giây. Khi đó $h=44{,}1$ mét.}]{Một quả bóng được đá lên từ độ cao $0$ m với vận tốc ban đầu theo phương thẳng đứng là $29{,}4$ m/s. Độ cao của bóng sau $t$ giây là $h(t)=-4{,}9t^2+29{,}4t+0$ (mét). Độ cao lớn nhất của quả bóng là bao nhiêu mét (làm tròn đến hàng phần trăm)?}
  {$44{,}1$.}
  {$45{,}3$.}{$43{,}3$.}{$29{,}4$.}

\cauTN[id=T10-C6-B16-D72-TH-012,muc=TH,loigiai={$a=-4<0$ nên $L$ lớn nhất tại $x=\dfrac{51}{2\cdot4}=\dfrac{51}{8}$, khi đó $L=672{,}56$ nghìn đồng.}]{Một cửa hàng ước tính lợi nhuận trong ngày thứ $x$ của tháng cho bởi $L(x)=-4x^2+51x+510$ (nghìn đồng), với $1\le x\le 30$. Lợi nhuận lớn nhất trong tháng là bao nhiêu nghìn đồng?}
  {$672{,}56$.}
  {$722{,}56$.}{$510$.}{$642{,}56$.}

\cauTN[id=T10-C6-B16-D72-TH-013,muc=TH,loigiai={Đồ thị là parabol có $a=-4{,}9<0$ nên $h$ đạt giá trị lớn nhất tại $t=-\dfrac{b}{2a}=\dfrac{29{,}4}{9{,}8}=3$ giây. Khi đó $h=44{,}1$ mét.}]{Một quả bóng được đá lên từ độ cao $0$ m với vận tốc ban đầu theo phương thẳng đứng là $29{,}4$ m/s. Độ cao của bóng sau $t$ giây là $h(t)=-4{,}9t^2+29{,}4t+0$ (mét). Độ cao lớn nhất của quả bóng là bao nhiêu mét (làm tròn đến hàng phần trăm)?}
  {$44{,}1$.}
  {$45{,}3$.}{$43{,}3$.}{$29{,}4$.}

\cauTN[id=T10-C6-B16-D72-TH-014,muc=TH,loigiai={$a=-2<0$ nên $L$ lớn nhất tại $x=\dfrac{72}{2\cdot2}=18$, khi đó $L=1\,402$ nghìn đồng.}]{Một cửa hàng ước tính lợi nhuận trong ngày thứ $x$ của tháng cho bởi $L(x)=-2x^2+72x+754$ (nghìn đồng), với $1\le x\le 30$. Lợi nhuận lớn nhất trong tháng là bao nhiêu nghìn đồng?}
  {$1\,402$.}
  {$1\,452$.}{$754$.}{$1\,372$.}

\cauTN[id=T10-C6-B16-D72-TH-015,muc=TH,loigiai={Đồ thị là parabol có $a=-4{,}9<0$ nên $h$ đạt giá trị lớn nhất tại $t=-\dfrac{b}{2a}=\dfrac{29{,}4}{9{,}8}=3$ giây. Khi đó $h=45{,}1$ mét.}]{Một quả bóng được đá lên từ độ cao $1$ m với vận tốc ban đầu theo phương thẳng đứng là $29{,}4$ m/s. Độ cao của bóng sau $t$ giây là $h(t)=-4{,}9t^2+29{,}4t+1$ (mét). Độ cao lớn nhất của quả bóng là bao nhiêu mét (làm tròn đến hàng phần trăm)?}
  {$45{,}1$.}
  {$46{,}3$.}{$44{,}3$.}{$29{,}4$.}

\cauTN[id=T10-C6-B16-D72-TH-016,muc=TH,loigiai={$a=-1<0$ nên $L$ lớn nhất tại $x=\dfrac{59}{2\cdot1}=\dfrac{59}{2}$, khi đó $L=1644{,}25$ nghìn đồng.}]{Một cửa hàng ước tính lợi nhuận trong ngày thứ $x$ của tháng cho bởi $L(x)=-x^2+59x+774$ (nghìn đồng), với $1\le x\le 30$. Lợi nhuận lớn nhất trong tháng là bao nhiêu nghìn đồng?}
  {$1644{,}25$.}
  {$1694{,}25$.}{$774$.}{$1614{,}25$.}

\cauTN[id=T10-C6-B16-D72-TH-017,muc=TH,loigiai={Đồ thị là parabol có $a=-4{,}9<0$ nên $h$ đạt giá trị lớn nhất tại $t=-\dfrac{b}{2a}=\dfrac{29{,}4}{9{,}8}=3$ giây. Khi đó $h=46{,}1$ mét.}]{Một quả bóng được đá lên từ độ cao $2$ m với vận tốc ban đầu theo phương thẳng đứng là $29{,}4$ m/s. Độ cao của bóng sau $t$ giây là $h(t)=-4{,}9t^2+29{,}4t+2$ (mét). Độ cao lớn nhất của quả bóng là bao nhiêu mét (làm tròn đến hàng phần trăm)?}
  {$46{,}1$.}
  {$47{,}3$.}{$45{,}3$.}{$29{,}4$.}

\cauTN[id=T10-C6-B16-D72-TH-018,muc=TH,loigiai={$a=-4<0$ nên $L$ lớn nhất tại $x=\dfrac{46}{2\cdot4}=\dfrac{23}{4}$, khi đó $L=655{,}25$ nghìn đồng.}]{Một cửa hàng ước tính lợi nhuận trong ngày thứ $x$ của tháng cho bởi $L(x)=-4x^2+46x+523$ (nghìn đồng), với $1\le x\le 30$. Lợi nhuận lớn nhất trong tháng là bao nhiêu nghìn đồng?}
  {$655{,}25$.}
  {$705{,}25$.}{$523$.}{$625{,}25$.}

\cauTN[id=T10-C6-B16-D72-TH-019,muc=TH,loigiai={Đồ thị là parabol có $a=-4{,}9<0$ nên $h$ đạt giá trị lớn nhất tại $t=-\dfrac{b}{2a}=\dfrac{24{,}5}{9{,}8}=2{,}5$ giây. Khi đó $h=31{,}62$ mét.}]{Một quả bóng được đá lên từ độ cao $1$ m với vận tốc ban đầu theo phương thẳng đứng là $24{,}5$ m/s. Độ cao của bóng sau $t$ giây là $h(t)=-4{,}9t^2+24{,}5t+1$ (mét). Độ cao lớn nhất của quả bóng là bao nhiêu mét (làm tròn đến hàng phần trăm)?}
  {$31{,}62$.}
  {$32{,}83$.}{$30{,}82$.}{$24{,}5$.}

\cauTN[id=T10-C6-B16-D72-TH-020,muc=TH,loigiai={$a=-3<0$ nên $L$ lớn nhất tại $x=\dfrac{44}{2\cdot3}=\dfrac{22}{3}$, khi đó $L=674{,}33$ nghìn đồng.}]{Một cửa hàng ước tính lợi nhuận trong ngày thứ $x$ của tháng cho bởi $L(x)=-3x^2+44x+513$ (nghìn đồng), với $1\le x\le 30$. Lợi nhuận lớn nhất trong tháng là bao nhiêu nghìn đồng?}
  {$674{,}33$.}
  {$724{,}33$.}{$513$.}{$644{,}33$.}

\cauTN[id=T10-C6-B16-D72-TH-021,muc=TH,loigiai={Đồ thị là parabol có $a=-4{,}9<0$ nên $h$ đạt giá trị lớn nhất tại $t=-\dfrac{b}{2a}=\dfrac{14{,}7}{9{,}8}=1{,}5$ giây. Khi đó $h=12{,}02$ mét.}]{Một quả bóng được đá lên từ độ cao $1$ m với vận tốc ban đầu theo phương thẳng đứng là $14{,}7$ m/s. Độ cao của bóng sau $t$ giây là $h(t)=-4{,}9t^2+14{,}7t+1$ (mét). Độ cao lớn nhất của quả bóng là bao nhiêu mét (làm tròn đến hàng phần trăm)?}
  {$12{,}02$.}
  {$13{,}22$.}{$11{,}22$.}{$14{,}7$.}

\cauTN[id=T10-C6-B16-D72-TH-022,muc=TH,loigiai={$a=-3<0$ nên $L$ lớn nhất tại $x=\dfrac{37}{2\cdot3}=\dfrac{37}{6}$, khi đó $L=373{,}08$ nghìn đồng.}]{Một cửa hàng ước tính lợi nhuận trong ngày thứ $x$ của tháng cho bởi $L(x)=-3x^2+37x+259$ (nghìn đồng), với $1\le x\le 30$. Lợi nhuận lớn nhất trong tháng là bao nhiêu nghìn đồng?}
  {$373{,}08$.}
  {$423{,}08$.}{$259$.}{$343{,}08$.}

\cauTN[id=T10-C6-B16-D72-TH-023,muc=TH,loigiai={Đồ thị là parabol có $a=-4{,}9<0$ nên $h$ đạt giá trị lớn nhất tại $t=-\dfrac{b}{2a}=\dfrac{24{,}5}{9{,}8}=2{,}5$ giây. Khi đó $h=30{,}62$ mét.}]{Một quả bóng được đá lên từ độ cao $0$ m với vận tốc ban đầu theo phương thẳng đứng là $24{,}5$ m/s. Độ cao của bóng sau $t$ giây là $h(t)=-4{,}9t^2+24{,}5t+0$ (mét). Độ cao lớn nhất của quả bóng là bao nhiêu mét (làm tròn đến hàng phần trăm)?}
  {$30{,}62$.}
  {$31{,}82$.}{$29{,}82$.}{$24{,}5$.}

\cauTN[id=T10-C6-B16-D72-TH-024,muc=TH,loigiai={$a=-3<0$ nên $L$ lớn nhất tại $x=\dfrac{48}{2\cdot3}=8$, khi đó $L=633$ nghìn đồng.}]{Một cửa hàng ước tính lợi nhuận trong ngày thứ $x$ của tháng cho bởi $L(x)=-3x^2+48x+441$ (nghìn đồng), với $1\le x\le 30$. Lợi nhuận lớn nhất trong tháng là bao nhiêu nghìn đồng?}
  {$633$.}
  {$683$.}{$441$.}{$603$.}

\cauDS[id=T10-C6-B16-D72-TH-025,muc=TH,loigiai={Đồ thị là parabol có $a=-4{,}9<0$ nên $h$ đạt giá trị lớn nhất tại $t=-\dfrac{b}{2a}=\dfrac{19{,}6}{9{,}8}=2$ giây. Khi đó $h=20{,}6$ mét.}]{Một quả bóng được đá lên từ độ cao $1$ m với vận tốc ban đầu theo phương thẳng đứng là $19{,}6$ m/s. Độ cao của bóng sau $t$ giây là $h(t)=-4{,}9t^2+19{,}6t+1$ (mét).}{%
  \y{Hệ số $a$ của hàm số là $-4{,}9$.}{d}
  \y{Bóng đạt độ cao lớn nhất tại thời điểm $t=2$ giây.}{d}
  \y{Độ cao ban đầu của bóng là $1$ mét.}{d}
  \y{Độ cao lớn nhất của bóng là $22{,}6$ mét.}{s}}

\cauDS[id=T10-C6-B16-D72-TH-026,muc=TH,loigiai={$a=-4<0$ nên $L$ lớn nhất tại $x=\dfrac{41}{2\cdot4}=\dfrac{41}{8}$, khi đó $L=428{,}06$ nghìn đồng.}]{Một cửa hàng ước tính lợi nhuận trong ngày thứ $x$ của tháng cho bởi $L(x)=-4x^2+41x+323$ (nghìn đồng), với $1\le x\le 30$.}{%
  \y{Lợi nhuận ngày đầu tháng là $360$ nghìn đồng.}{d}
  \y{Lợi nhuận lớn nhất đạt được ở ngày thứ $\dfrac{41}{8}$.}{d}
  \y{Hàm số $L(x)$ đồng biến trên toàn bộ đoạn $[1;30]$.}{s}
  \y{Lợi nhuận lớn nhất là $428{,}06$ nghìn đồng.}{d}}

\cauDS[id=T10-C6-B16-D72-TH-027,muc=TH,loigiai={Đồ thị là parabol có $a=-4{,}9<0$ nên $h$ đạt giá trị lớn nhất tại $t=-\dfrac{b}{2a}=\dfrac{24{,}5}{9{,}8}=2{,}5$ giây. Khi đó $h=32{,}62$ mét.}]{Một quả bóng được đá lên từ độ cao $2$ m với vận tốc ban đầu theo phương thẳng đứng là $24{,}5$ m/s. Độ cao của bóng sau $t$ giây là $h(t)=-4{,}9t^2+24{,}5t+2$ (mét).}{%
  \y{Hệ số $a$ của hàm số là $-4{,}9$.}{d}
  \y{Bóng đạt độ cao lớn nhất tại thời điểm $t=2{,}5$ giây.}{d}
  \y{Độ cao ban đầu của bóng là $2$ mét.}{d}
  \y{Độ cao lớn nhất của bóng là $34{,}62$ mét.}{s}}

\cauDS[id=T10-C6-B16-D72-TH-028,muc=TH,loigiai={$a=-4<0$ nên $L$ lớn nhất tại $x=\dfrac{55}{2\cdot4}=\dfrac{55}{8}$, khi đó $L=863{,}06$ nghìn đồng.}]{Một cửa hàng ước tính lợi nhuận trong ngày thứ $x$ của tháng cho bởi $L(x)=-4x^2+55x+674$ (nghìn đồng), với $1\le x\le 30$.}{%
  \y{Lợi nhuận ngày đầu tháng là $725$ nghìn đồng.}{d}
  \y{Lợi nhuận lớn nhất đạt được ở ngày thứ $\dfrac{55}{8}$.}{d}
  \y{Hàm số $L(x)$ đồng biến trên toàn bộ đoạn $[1;30]$.}{s}
  \y{Lợi nhuận lớn nhất là $863{,}06$ nghìn đồng.}{d}}

\cauDS[id=T10-C6-B16-D72-TH-029,muc=TH,loigiai={Đồ thị là parabol có $a=-4{,}9<0$ nên $h$ đạt giá trị lớn nhất tại $t=-\dfrac{b}{2a}=\dfrac{14{,}7}{9{,}8}=1{,}5$ giây. Khi đó $h=12{,}02$ mét.}]{Một quả bóng được đá lên từ độ cao $1$ m với vận tốc ban đầu theo phương thẳng đứng là $14{,}7$ m/s. Độ cao của bóng sau $t$ giây là $h(t)=-4{,}9t^2+14{,}7t+1$ (mét).}{%
  \y{Hệ số $a$ của hàm số là $-4{,}9$.}{d}
  \y{Bóng đạt độ cao lớn nhất tại thời điểm $t=1{,}5$ giây.}{d}
  \y{Độ cao ban đầu của bóng là $1$ mét.}{d}
  \y{Độ cao lớn nhất của bóng là $14{,}02$ mét.}{s}}

\cauDS[id=T10-C6-B16-D72-TH-030,muc=TH,loigiai={$a=-2<0$ nên $L$ lớn nhất tại $x=\dfrac{69}{2\cdot2}=\dfrac{69}{4}$, khi đó $L=1451{,}12$ nghìn đồng.}]{Một cửa hàng ước tính lợi nhuận trong ngày thứ $x$ của tháng cho bởi $L(x)=-2x^2+69x+856$ (nghìn đồng), với $1\le x\le 30$.}{%
  \y{Lợi nhuận ngày đầu tháng là $923$ nghìn đồng.}{d}
  \y{Lợi nhuận lớn nhất đạt được ở ngày thứ $\dfrac{69}{4}$.}{d}
  \y{Hàm số $L(x)$ đồng biến trên toàn bộ đoạn $[1;30]$.}{s}
  \y{Lợi nhuận lớn nhất là $1451{,}12$ nghìn đồng.}{d}}

\cauDS[id=T10-C6-B16-D72-TH-031,muc=TH,loigiai={Đồ thị là parabol có $a=-4{,}9<0$ nên $h$ đạt giá trị lớn nhất tại $t=-\dfrac{b}{2a}=\dfrac{14{,}7}{9{,}8}=1{,}5$ giây. Khi đó $h=12{,}52$ mét.}]{Một quả bóng được đá lên từ độ cao $1{,}5$ m với vận tốc ban đầu theo phương thẳng đứng là $14{,}7$ m/s. Độ cao của bóng sau $t$ giây là $h(t)=-4{,}9t^2+14{,}7t+1{,}5$ (mét).}{%
  \y{Hệ số $a$ của hàm số là $-4{,}9$.}{d}
  \y{Bóng đạt độ cao lớn nhất tại thời điểm $t=1{,}5$ giây.}{d}
  \y{Độ cao ban đầu của bóng là $1{,}5$ mét.}{d}
  \y{Độ cao lớn nhất của bóng là $14{,}52$ mét.}{s}}

\cauDS[id=T10-C6-B16-D72-TH-032,muc=TH,loigiai={$a=-1<0$ nên $L$ lớn nhất tại $x=\dfrac{56}{2\cdot1}=28$, khi đó $L=1\,048$ nghìn đồng.}]{Một cửa hàng ước tính lợi nhuận trong ngày thứ $x$ của tháng cho bởi $L(x)=-x^2+56x+264$ (nghìn đồng), với $1\le x\le 30$.}{%
  \y{Lợi nhuận ngày đầu tháng là $319$ nghìn đồng.}{d}
  \y{Lợi nhuận lớn nhất đạt được ở ngày thứ $28$.}{d}
  \y{Hàm số $L(x)$ đồng biến trên toàn bộ đoạn $[1;30]$.}{s}
  \y{Lợi nhuận lớn nhất là $1\,048$ nghìn đồng.}{d}}

\cauTLN[id=T10-C6-B16-D72-TH-033,muc=TH,loigiai={Đồ thị là parabol có $a=-4{,}9<0$ nên $h$ đạt giá trị lớn nhất tại $t=-\dfrac{b}{2a}=\dfrac{24{,}5}{9{,}8}=2{,}5$ giây. Khi đó $h=32{,}12$ mét.}]{Một quả bóng được đá lên từ độ cao $1{,}5$ m với vận tốc ban đầu theo phương thẳng đứng là $24{,}5$ m/s. Độ cao của bóng sau $t$ giây là $h(t)=-4{,}9t^2+24{,}5t+1{,}5$ (mét). Độ cao lớn nhất của quả bóng là bao nhiêu mét (làm tròn đến hàng phần trăm)?}{32,12}

\cauTLN[id=T10-C6-B16-D72-TH-034,muc=TH,loigiai={$a=-4<0$ nên $L$ lớn nhất tại $x=\dfrac{51}{2\cdot4}=\dfrac{51}{8}$, khi đó $L=405{,}56$ nghìn đồng.}]{Một cửa hàng ước tính lợi nhuận trong ngày thứ $x$ của tháng cho bởi $L(x)=-4x^2+51x+243$ (nghìn đồng), với $1\le x\le 30$. Lợi nhuận lớn nhất trong tháng là bao nhiêu nghìn đồng?}{405,56}

\cauTLN[id=T10-C6-B16-D72-TH-035,muc=TH,loigiai={Đồ thị là parabol có $a=-4{,}9<0$ nên $h$ đạt giá trị lớn nhất tại $t=-\dfrac{b}{2a}=\dfrac{14{,}7}{9{,}8}=1{,}5$ giây. Khi đó $h=12{,}02$ mét.}]{Một quả bóng được đá lên từ độ cao $1$ m với vận tốc ban đầu theo phương thẳng đứng là $14{,}7$ m/s. Độ cao của bóng sau $t$ giây là $h(t)=-4{,}9t^2+14{,}7t+1$ (mét). Độ cao lớn nhất của quả bóng là bao nhiêu mét (làm tròn đến hàng phần trăm)?}{12,02}

\cauTLN[id=T10-C6-B16-D72-TH-036,muc=TH,loigiai={$a=-4<0$ nên $L$ lớn nhất tại $x=\dfrac{39}{2\cdot4}=\dfrac{39}{8}$, khi đó $L=616{,}06$ nghìn đồng.}]{Một cửa hàng ước tính lợi nhuận trong ngày thứ $x$ của tháng cho bởi $L(x)=-4x^2+39x+521$ (nghìn đồng), với $1\le x\le 30$. Lợi nhuận lớn nhất trong tháng là bao nhiêu nghìn đồng?}{616,06}

\cauTLN[id=T10-C6-B16-D72-TH-037,muc=TH,loigiai={Đồ thị là parabol có $a=-4{,}9<0$ nên $h$ đạt giá trị lớn nhất tại $t=-\dfrac{b}{2a}=\dfrac{14{,}7}{9{,}8}=1{,}5$ giây. Khi đó $h=12{,}02$ mét.}]{Một quả bóng được đá lên từ độ cao $1$ m với vận tốc ban đầu theo phương thẳng đứng là $14{,}7$ m/s. Độ cao của bóng sau $t$ giây là $h(t)=-4{,}9t^2+14{,}7t+1$ (mét). Độ cao lớn nhất của quả bóng là bao nhiêu mét (làm tròn đến hàng phần trăm)?}{12,02}

\cauTLN[id=T10-C6-B16-D72-TH-038,muc=TH,loigiai={$a=-2<0$ nên $L$ lớn nhất tại $x=\dfrac{52}{2\cdot2}=13$, khi đó $L=1\,064$ nghìn đồng.}]{Một cửa hàng ước tính lợi nhuận trong ngày thứ $x$ của tháng cho bởi $L(x)=-2x^2+52x+726$ (nghìn đồng), với $1\le x\le 30$. Lợi nhuận lớn nhất trong tháng là bao nhiêu nghìn đồng?}{1064}

\cauTLN[id=T10-C6-B16-D72-TH-039,muc=TH,loigiai={Đồ thị là parabol có $a=-4{,}9<0$ nên $h$ đạt giá trị lớn nhất tại $t=-\dfrac{b}{2a}=\dfrac{24{,}5}{9{,}8}=2{,}5$ giây. Khi đó $h=32{,}62$ mét.}]{Một quả bóng được đá lên từ độ cao $2$ m với vận tốc ban đầu theo phương thẳng đứng là $24{,}5$ m/s. Độ cao của bóng sau $t$ giây là $h(t)=-4{,}9t^2+24{,}5t+2$ (mét). Độ cao lớn nhất của quả bóng là bao nhiêu mét (làm tròn đến hàng phần trăm)?}{32,62}

\cauTLN[id=T10-C6-B16-D72-TH-040,muc=TH,loigiai={$a=-4<0$ nên $L$ lớn nhất tại $x=\dfrac{69}{2\cdot4}=\dfrac{69}{8}$, khi đó $L=961{,}56$ nghìn đồng.}]{Một cửa hàng ước tính lợi nhuận trong ngày thứ $x$ của tháng cho bởi $L(x)=-4x^2+69x+664$ (nghìn đồng), với $1\le x\le 30$. Lợi nhuận lớn nhất trong tháng là bao nhiêu nghìn đồng?}{961,56}

\cauTN[id=T10-C6-B16-D72-VD-001,muc=VD,loigiai={$a=-2<0$ nên $L$ lớn nhất tại $x=\dfrac{82}{2\cdot2}=\dfrac{41}{2}$, khi đó $L=1464{,}5$ nghìn đồng.}]{Một cửa hàng ước tính lợi nhuận trong ngày thứ $x$ của tháng cho bởi $L(x)=-2x^2+82x+624$ (nghìn đồng), với $1\le x\le 30$. Lợi nhuận lớn nhất trong tháng là bao nhiêu nghìn đồng?}
  {$1464{,}5$.}
  {$1514{,}5$.}{$624$.}{$1434{,}5$.}

\cauTN[id=T10-C6-B16-D72-VD-002,muc=VD,loigiai={Đồ thị là parabol có $a=-4{,}9<0$ nên $h$ đạt giá trị lớn nhất tại $t=-\dfrac{b}{2a}=\dfrac{24{,}5}{9{,}8}=2{,}5$ giây. Khi đó $h=30{,}62$ mét.}]{Một quả bóng được đá lên từ độ cao $0$ m với vận tốc ban đầu theo phương thẳng đứng là $24{,}5$ m/s. Độ cao của bóng sau $t$ giây là $h(t)=-4{,}9t^2+24{,}5t+0$ (mét). Độ cao lớn nhất của quả bóng là bao nhiêu mét (làm tròn đến hàng phần trăm)?}
  {$30{,}62$.}
  {$31{,}82$.}{$29{,}82$.}{$24{,}5$.}

\cauTN[id=T10-C6-B16-D72-VD-003,muc=VD,loigiai={$a=-1<0$ nên $L$ lớn nhất tại $x=\dfrac{42}{2\cdot1}=21$, khi đó $L=1\,126$ nghìn đồng.}]{Một cửa hàng ước tính lợi nhuận trong ngày thứ $x$ của tháng cho bởi $L(x)=-x^2+42x+685$ (nghìn đồng), với $1\le x\le 30$. Lợi nhuận lớn nhất trong tháng là bao nhiêu nghìn đồng?}
  {$1\,126$.}
  {$1\,176$.}{$685$.}{$1\,096$.}

\cauTN[id=T10-C6-B16-D72-VD-004,muc=VD,loigiai={Đồ thị là parabol có $a=-4{,}9<0$ nên $h$ đạt giá trị lớn nhất tại $t=-\dfrac{b}{2a}=\dfrac{14{,}7}{9{,}8}=1{,}5$ giây. Khi đó $h=12{,}02$ mét.}]{Một quả bóng được đá lên từ độ cao $1$ m với vận tốc ban đầu theo phương thẳng đứng là $14{,}7$ m/s. Độ cao của bóng sau $t$ giây là $h(t)=-4{,}9t^2+14{,}7t+1$ (mét). Độ cao lớn nhất của quả bóng là bao nhiêu mét (làm tròn đến hàng phần trăm)?}
  {$12{,}02$.}
  {$13{,}22$.}{$11{,}22$.}{$14{,}7$.}

\cauTN[id=T10-C6-B16-D72-VD-005,muc=VD,loigiai={$a=-1<0$ nên $L$ lớn nhất tại $x=\dfrac{36}{2\cdot1}=18$, khi đó $L=1\,171$ nghìn đồng.}]{Một cửa hàng ước tính lợi nhuận trong ngày thứ $x$ của tháng cho bởi $L(x)=-x^2+36x+847$ (nghìn đồng), với $1\le x\le 30$. Lợi nhuận lớn nhất trong tháng là bao nhiêu nghìn đồng?}
  {$1\,171$.}
  {$1\,221$.}{$847$.}{$1\,141$.}

\cauTN[id=T10-C6-B16-D72-VD-006,muc=VD,loigiai={Đồ thị là parabol có $a=-4{,}9<0$ nên $h$ đạt giá trị lớn nhất tại $t=-\dfrac{b}{2a}=\dfrac{14{,}7}{9{,}8}=1{,}5$ giây. Khi đó $h=11{,}02$ mét.}]{Một quả bóng được đá lên từ độ cao $0$ m với vận tốc ban đầu theo phương thẳng đứng là $14{,}7$ m/s. Độ cao của bóng sau $t$ giây là $h(t)=-4{,}9t^2+14{,}7t+0$ (mét). Độ cao lớn nhất của quả bóng là bao nhiêu mét (làm tròn đến hàng phần trăm)?}
  {$11{,}02$.}
  {$12{,}22$.}{$10{,}22$.}{$14{,}7$.}

\cauTN[id=T10-C6-B16-D72-VD-007,muc=VD,loigiai={$a=-4<0$ nên $L$ lớn nhất tại $x=\dfrac{34}{2\cdot4}=\dfrac{17}{4}$, khi đó $L=586{,}25$ nghìn đồng.}]{Một cửa hàng ước tính lợi nhuận trong ngày thứ $x$ của tháng cho bởi $L(x)=-4x^2+34x+514$ (nghìn đồng), với $1\le x\le 30$. Lợi nhuận lớn nhất trong tháng là bao nhiêu nghìn đồng?}
  {$586{,}25$.}
  {$636{,}25$.}{$514$.}{$556{,}25$.}

\cauTN[id=T10-C6-B16-D72-VD-008,muc=VD,loigiai={Đồ thị là parabol có $a=-4{,}9<0$ nên $h$ đạt giá trị lớn nhất tại $t=-\dfrac{b}{2a}=\dfrac{14{,}7}{9{,}8}=1{,}5$ giây. Khi đó $h=12{,}02$ mét.}]{Một quả bóng được đá lên từ độ cao $1$ m với vận tốc ban đầu theo phương thẳng đứng là $14{,}7$ m/s. Độ cao của bóng sau $t$ giây là $h(t)=-4{,}9t^2+14{,}7t+1$ (mét). Độ cao lớn nhất của quả bóng là bao nhiêu mét (làm tròn đến hàng phần trăm)?}
  {$12{,}02$.}
  {$13{,}22$.}{$11{,}22$.}{$14{,}7$.}

\cauTN[id=T10-C6-B16-D72-VD-009,muc=VD,loigiai={$a=-4<0$ nên $L$ lớn nhất tại $x=\dfrac{45}{2\cdot4}=\dfrac{45}{8}$, khi đó $L=330{,}56$ nghìn đồng.}]{Một cửa hàng ước tính lợi nhuận trong ngày thứ $x$ của tháng cho bởi $L(x)=-4x^2+45x+204$ (nghìn đồng), với $1\le x\le 30$. Lợi nhuận lớn nhất trong tháng là bao nhiêu nghìn đồng?}
  {$330{,}56$.}
  {$380{,}56$.}{$204$.}{$300{,}56$.}

\cauTN[id=T10-C6-B16-D72-VD-010,muc=VD,loigiai={Đồ thị là parabol có $a=-4{,}9<0$ nên $h$ đạt giá trị lớn nhất tại $t=-\dfrac{b}{2a}=\dfrac{29{,}4}{9{,}8}=3$ giây. Khi đó $h=45{,}1$ mét.}]{Một quả bóng được đá lên từ độ cao $1$ m với vận tốc ban đầu theo phương thẳng đứng là $29{,}4$ m/s. Độ cao của bóng sau $t$ giây là $h(t)=-4{,}9t^2+29{,}4t+1$ (mét). Độ cao lớn nhất của quả bóng là bao nhiêu mét (làm tròn đến hàng phần trăm)?}
  {$45{,}1$.}
  {$46{,}3$.}{$44{,}3$.}{$29{,}4$.}

\cauTN[id=T10-C6-B16-D72-VD-011,muc=VD,loigiai={$a=-3<0$ nên $L$ lớn nhất tại $x=\dfrac{68}{2\cdot3}=\dfrac{34}{3}$, khi đó $L=1240{,}33$ nghìn đồng.}]{Một cửa hàng ước tính lợi nhuận trong ngày thứ $x$ của tháng cho bởi $L(x)=-3x^2+68x+855$ (nghìn đồng), với $1\le x\le 30$. Lợi nhuận lớn nhất trong tháng là bao nhiêu nghìn đồng?}
  {$1240{,}33$.}
  {$1290{,}33$.}{$855$.}{$1210{,}33$.}

\cauTN[id=T10-C6-B16-D72-VD-012,muc=VD,loigiai={Đồ thị là parabol có $a=-4{,}9<0$ nên $h$ đạt giá trị lớn nhất tại $t=-\dfrac{b}{2a}=\dfrac{14{,}7}{9{,}8}=1{,}5$ giây. Khi đó $h=12{,}02$ mét.}]{Một quả bóng được đá lên từ độ cao $1$ m với vận tốc ban đầu theo phương thẳng đứng là $14{,}7$ m/s. Độ cao của bóng sau $t$ giây là $h(t)=-4{,}9t^2+14{,}7t+1$ (mét). Độ cao lớn nhất của quả bóng là bao nhiêu mét (làm tròn đến hàng phần trăm)?}
  {$12{,}02$.}
  {$13{,}22$.}{$11{,}22$.}{$14{,}7$.}

\cauDS[id=T10-C6-B16-D72-VD-013,muc=VD,loigiai={$a=-4<0$ nên $L$ lớn nhất tại $x=\dfrac{89}{2\cdot4}=\dfrac{89}{8}$, khi đó $L=1331{,}06$ nghìn đồng.}]{Một cửa hàng ước tính lợi nhuận trong ngày thứ $x$ của tháng cho bởi $L(x)=-4x^2+89x+836$ (nghìn đồng), với $1\le x\le 30$.}{%
  \y{Lợi nhuận ngày đầu tháng là $921$ nghìn đồng.}{d}
  \y{Lợi nhuận lớn nhất đạt được ở ngày thứ $\dfrac{89}{8}$.}{d}
  \y{Hàm số $L(x)$ đồng biến trên toàn bộ đoạn $[1;30]$.}{s}
  \y{Lợi nhuận lớn nhất là $1331{,}06$ nghìn đồng.}{d}}

\cauDS[id=T10-C6-B16-D72-VD-014,muc=VD,loigiai={Đồ thị là parabol có $a=-4{,}9<0$ nên $h$ đạt giá trị lớn nhất tại $t=-\dfrac{b}{2a}=\dfrac{24{,}5}{9{,}8}=2{,}5$ giây. Khi đó $h=32{,}12$ mét.}]{Một quả bóng được đá lên từ độ cao $1{,}5$ m với vận tốc ban đầu theo phương thẳng đứng là $24{,}5$ m/s. Độ cao của bóng sau $t$ giây là $h(t)=-4{,}9t^2+24{,}5t+1{,}5$ (mét).}{%
  \y{Hệ số $a$ của hàm số là $-4{,}9$.}{d}
  \y{Bóng đạt độ cao lớn nhất tại thời điểm $t=2{,}5$ giây.}{d}
  \y{Độ cao ban đầu của bóng là $1{,}5$ mét.}{d}
  \y{Độ cao lớn nhất của bóng là $34{,}12$ mét.}{s}}

\cauDS[id=T10-C6-B16-D72-VD-015,muc=VD,loigiai={$a=-4<0$ nên $L$ lớn nhất tại $x=\dfrac{72}{2\cdot4}=9$, khi đó $L=730$ nghìn đồng.}]{Một cửa hàng ước tính lợi nhuận trong ngày thứ $x$ của tháng cho bởi $L(x)=-4x^2+72x+406$ (nghìn đồng), với $1\le x\le 30$.}{%
  \y{Lợi nhuận ngày đầu tháng là $474$ nghìn đồng.}{d}
  \y{Lợi nhuận lớn nhất đạt được ở ngày thứ $9$.}{d}
  \y{Hàm số $L(x)$ đồng biến trên toàn bộ đoạn $[1;30]$.}{s}
  \y{Lợi nhuận lớn nhất là $730$ nghìn đồng.}{d}}

\cauDS[id=T10-C6-B16-D72-VD-016,muc=VD,loigiai={Đồ thị là parabol có $a=-4{,}9<0$ nên $h$ đạt giá trị lớn nhất tại $t=-\dfrac{b}{2a}=\dfrac{24{,}5}{9{,}8}=2{,}5$ giây. Khi đó $h=30{,}62$ mét.}]{Một quả bóng được đá lên từ độ cao $0$ m với vận tốc ban đầu theo phương thẳng đứng là $24{,}5$ m/s. Độ cao của bóng sau $t$ giây là $h(t)=-4{,}9t^2+24{,}5t+0$ (mét).}{%
  \y{Hệ số $a$ của hàm số là $-4{,}9$.}{d}
  \y{Bóng đạt độ cao lớn nhất tại thời điểm $t=2{,}5$ giây.}{d}
  \y{Độ cao ban đầu của bóng là $0$ mét.}{d}
  \y{Độ cao lớn nhất của bóng là $32{,}62$ mét.}{s}}

\cauTLN[id=T10-C6-B16-D72-VD-017,muc=VD,loigiai={$a=-1<0$ nên $L$ lớn nhất tại $x=\dfrac{80}{2\cdot1}=40$, khi đó $L=2\,268$ nghìn đồng.}]{Một cửa hàng ước tính lợi nhuận trong ngày thứ $x$ của tháng cho bởi $L(x)=-x^2+80x+668$ (nghìn đồng), với $1\le x\le 30$. Lợi nhuận lớn nhất trong tháng là bao nhiêu nghìn đồng?}{2268}

\cauTLN[id=T10-C6-B16-D72-VD-018,muc=VD,loigiai={Đồ thị là parabol có $a=-4{,}9<0$ nên $h$ đạt giá trị lớn nhất tại $t=-\dfrac{b}{2a}=\dfrac{14{,}7}{9{,}8}=1{,}5$ giây. Khi đó $h=11{,}02$ mét.}]{Một quả bóng được đá lên từ độ cao $0$ m với vận tốc ban đầu theo phương thẳng đứng là $14{,}7$ m/s. Độ cao của bóng sau $t$ giây là $h(t)=-4{,}9t^2+14{,}7t+0$ (mét). Độ cao lớn nhất của quả bóng là bao nhiêu mét (làm tròn đến hàng phần trăm)?}{11,02}

\cauTLN[id=T10-C6-B16-D72-VD-019,muc=VD,loigiai={$a=-2<0$ nên $L$ lớn nhất tại $x=\dfrac{44}{2\cdot2}=11$, khi đó $L=799$ nghìn đồng.}]{Một cửa hàng ước tính lợi nhuận trong ngày thứ $x$ của tháng cho bởi $L(x)=-2x^2+44x+557$ (nghìn đồng), với $1\le x\le 30$. Lợi nhuận lớn nhất trong tháng là bao nhiêu nghìn đồng?}{799}

\cauTLN[id=T10-C6-B16-D72-VD-020,muc=VD,loigiai={Đồ thị là parabol có $a=-4{,}9<0$ nên $h$ đạt giá trị lớn nhất tại $t=-\dfrac{b}{2a}=\dfrac{19{,}6}{9{,}8}=2$ giây. Khi đó $h=20{,}6$ mét.}]{Một quả bóng được đá lên từ độ cao $1$ m với vận tốc ban đầu theo phương thẳng đứng là $19{,}6$ m/s. Độ cao của bóng sau $t$ giây là $h(t)=-4{,}9t^2+19{,}6t+1$ (mét). Độ cao lớn nhất của quả bóng là bao nhiêu mét (làm tròn đến hàng phần trăm)?}{20,6}

\cauTN[id=T10-C6-B16-D72-VDC-001,muc=VDC,loigiai={$a=-1<0$ nên $L$ lớn nhất tại $x=\dfrac{69}{2\cdot1}=\dfrac{69}{2}$, khi đó $L=1459{,}25$ nghìn đồng.}]{Một cửa hàng ước tính lợi nhuận trong ngày thứ $x$ của tháng cho bởi $L(x)=-x^2+69x+269$ (nghìn đồng), với $1\le x\le 30$. Lợi nhuận lớn nhất trong tháng là bao nhiêu nghìn đồng?}
  {$1459{,}25$.}
  {$1509{,}25$.}{$269$.}{$1429{,}25$.}

\cauTN[id=T10-C6-B16-D72-VDC-002,muc=VDC,loigiai={$a=-3<0$ nên $L$ lớn nhất tại $x=\dfrac{75}{2\cdot3}=\dfrac{25}{2}$, khi đó $L=1094{,}75$ nghìn đồng.}]{Một cửa hàng ước tính lợi nhuận trong ngày thứ $x$ của tháng cho bởi $L(x)=-3x^2+75x+626$ (nghìn đồng), với $1\le x\le 30$. Lợi nhuận lớn nhất trong tháng là bao nhiêu nghìn đồng?}
  {$1094{,}75$.}
  {$1144{,}75$.}{$626$.}{$1064{,}75$.}

\cauTN[id=T10-C6-B16-D72-VDC-003,muc=VDC,loigiai={$a=-1<0$ nên $L$ lớn nhất tại $x=\dfrac{84}{2\cdot1}=42$, khi đó $L=2\,166$ nghìn đồng.}]{Một cửa hàng ước tính lợi nhuận trong ngày thứ $x$ của tháng cho bởi $L(x)=-x^2+84x+402$ (nghìn đồng), với $1\le x\le 30$. Lợi nhuận lớn nhất trong tháng là bao nhiêu nghìn đồng?}
  {$2\,166$.}
  {$2\,216$.}{$402$.}{$2\,136$.}

\cauTN[id=T10-C6-B16-D72-VDC-004,muc=VDC,loigiai={$a=-1<0$ nên $L$ lớn nhất tại $x=\dfrac{38}{2\cdot1}=19$, khi đó $L=885$ nghìn đồng.}]{Một cửa hàng ước tính lợi nhuận trong ngày thứ $x$ của tháng cho bởi $L(x)=-x^2+38x+524$ (nghìn đồng), với $1\le x\le 30$. Lợi nhuận lớn nhất trong tháng là bao nhiêu nghìn đồng?}
  {$885$.}
  {$935$.}{$524$.}{$855$.}

\cauTN[id=T10-C6-B16-D72-VDC-005,muc=VDC,loigiai={$a=-3<0$ nên $L$ lớn nhất tại $x=\dfrac{74}{2\cdot3}=\dfrac{37}{3}$, khi đó $L=610{,}33$ nghìn đồng.}]{Một cửa hàng ước tính lợi nhuận trong ngày thứ $x$ của tháng cho bởi $L(x)=-3x^2+74x+154$ (nghìn đồng), với $1\le x\le 30$. Lợi nhuận lớn nhất trong tháng là bao nhiêu nghìn đồng?}
  {$610{,}33$.}
  {$660{,}33$.}{$154$.}{$580{,}33$.}

\cauTN[id=T10-C6-B16-D72-VDC-006,muc=VDC,loigiai={$a=-3<0$ nên $L$ lớn nhất tại $x=\dfrac{37}{2\cdot3}=\dfrac{37}{6}$, khi đó $L=573{,}08$ nghìn đồng.}]{Một cửa hàng ước tính lợi nhuận trong ngày thứ $x$ của tháng cho bởi $L(x)=-3x^2+37x+459$ (nghìn đồng), với $1\le x\le 30$. Lợi nhuận lớn nhất trong tháng là bao nhiêu nghìn đồng?}
  {$573{,}08$.}
  {$623{,}08$.}{$459$.}{$543{,}08$.}

\cauDS[id=T10-C6-B16-D72-VDC-007,muc=VDC,loigiai={$a=-1<0$ nên $L$ lớn nhất tại $x=\dfrac{40}{2\cdot1}=20$, khi đó $L=1\,154$ nghìn đồng.}]{Một cửa hàng ước tính lợi nhuận trong ngày thứ $x$ của tháng cho bởi $L(x)=-x^2+40x+754$ (nghìn đồng), với $1\le x\le 30$.}{%
  \y{Lợi nhuận ngày đầu tháng là $793$ nghìn đồng.}{d}
  \y{Lợi nhuận lớn nhất đạt được ở ngày thứ $20$.}{d}
  \y{Hàm số $L(x)$ đồng biến trên toàn bộ đoạn $[1;30]$.}{s}
  \y{Lợi nhuận lớn nhất là $1\,154$ nghìn đồng.}{d}}

\cauDS[id=T10-C6-B16-D72-VDC-008,muc=VDC,loigiai={$a=-4<0$ nên $L$ lớn nhất tại $x=\dfrac{69}{2\cdot4}=\dfrac{69}{8}$, khi đó $L=935{,}56$ nghìn đồng.}]{Một cửa hàng ước tính lợi nhuận trong ngày thứ $x$ của tháng cho bởi $L(x)=-4x^2+69x+638$ (nghìn đồng), với $1\le x\le 30$.}{%
  \y{Lợi nhuận ngày đầu tháng là $703$ nghìn đồng.}{d}
  \y{Lợi nhuận lớn nhất đạt được ở ngày thứ $\dfrac{69}{8}$.}{d}
  \y{Hàm số $L(x)$ đồng biến trên toàn bộ đoạn $[1;30]$.}{s}
  \y{Lợi nhuận lớn nhất là $935{,}56$ nghìn đồng.}{d}}

\cauTLN[id=T10-C6-B16-D72-VDC-009,muc=VDC,loigiai={$a=-3<0$ nên $L$ lớn nhất tại $x=\dfrac{44}{2\cdot3}=\dfrac{22}{3}$, khi đó $L=974{,}33$ nghìn đồng.}]{Một cửa hàng ước tính lợi nhuận trong ngày thứ $x$ của tháng cho bởi $L(x)=-3x^2+44x+813$ (nghìn đồng), với $1\le x\le 30$. Lợi nhuận lớn nhất trong tháng là bao nhiêu nghìn đồng?}{974,33}

\cauTLN[id=T10-C6-B16-D72-VDC-010,muc=VDC,loigiai={$a=-2<0$ nên $L$ lớn nhất tại $x=\dfrac{31}{2\cdot2}=\dfrac{31}{4}$, khi đó $L=712{,}12$ nghìn đồng.}]{Một cửa hàng ước tính lợi nhuận trong ngày thứ $x$ của tháng cho bởi $L(x)=-2x^2+31x+592$ (nghìn đồng), với $1\le x\le 30$. Lợi nhuận lớn nhất trong tháng là bao nhiêu nghìn đồng?}{712,12}
